\specialchapter{TÓM TẮT}

Người tìm đồ thời trang thường mô tả món đồ bằng màu sắc, dáng và phong cách, nên tìm theo từ khoá dễ trượt. Đồ án xây dựng hệ thống tìm ảnh sản phẩm thời trang bằng truy vấn tiếng Việt và đo xem biểu diễn chung ảnh và văn bản có tốt hơn BM25 hay không. Catalog là 6.963 ảnh kiểm tra của Fashionpedia. Bộ đánh giá gồm 155 truy vấn tiếng Việt do nhóm soạn, với nhãn do hai người chấm độc lập và một bộ lọc tự động gán. Hệ thống dùng SigLIP~2 để đặt truy vấn và ảnh vào cùng không gian vector, tinh chỉnh bằng adapter DoRA, và được đối chiếu với LoRA thường, CLIP, ViSigLIP-OT và BM25 chạy trên tên danh mục. SigLIP~2 dùng sẵn đạt Recall@10 là 0,865 so với 0,568 của BM25 ($p < 10^{-4}$), còn CLIP và ViSigLIP-OT không hơn BM25. Tinh chỉnh nâng Recall@1 từ 0,490 lên 0,697 ($p < 10^{-4}$); Recall@10 chỉ tăng từ 0,865 lên 0,884 ($p = 0{,}240$). LoRA thường đạt Recall@1 là 0,690, chưa phân biệt được với DoRA ($p = 0{,}82$), nên mức tăng chưa quy được riêng cho DoRA. Ba điểm cần đọc kèm: bộ lọc nhãn tự động làm lệch 35 truy vấn, trong đó 26 truy vấn bị chấm sai và chứa cả 18 truy vấn mà hệ DoRA trượt ở top-10; mỗi cấu hình chỉ huấn luyện một lần; các giá trị $p$ chưa hiệu chỉnh cho nhiều phép so sánh. Đóng góp gồm bộ đánh giá tiếng Việt cho tìm ảnh thời trang, so sánh có kiểm định giữa các hướng và đối soát nhãn tự động.

\medskip
\noindent\textbf{Từ khoá:} tìm kiếm đa mô thức, tìm ảnh thời trang, SigLIP~2, DoRA, truy vấn tiếng Việt, Fashionpedia.
